\documentclass[11pt]{article}
\usepackage[margin=1in]{geometry}
\usepackage{amsmath,amssymb,amsthm,mathtools}
\usepackage[T1]{fontenc}
\usepackage[utf8]{inputenc}
\usepackage{enumitem}
\usepackage{booktabs,longtable,array}
\usepackage{hyperref}
\usepackage{xcolor}
\usepackage{microtype}
\hypersetup{colorlinks=true, linkcolor=blue!50!black, citecolor=blue!50!black, urlcolor=blue!50!black}

\newtheorem{theorem}{Theorem}[section]
\newtheorem{proposition}[theorem]{Proposition}
\newtheorem{lemma}[theorem]{Lemma}
\newtheorem{corollary}[theorem]{Corollary}
\newtheorem{claim}[theorem]{Claim}
\theoremstyle{definition}
\newtheorem{definition}[theorem]{Definition}
\newtheorem{assumption}[theorem]{Assumption}
\theoremstyle{remark}
\newtheorem{remark}[theorem]{Remark}

\newcommand{\KL}{\mathrm{KL}}
\newcommand{\Ent}{\mathrm{Ent}}
\newcommand{\Prob}{\mathsf{P}}
\newcommand{\E}{\mathbb{E}}
\newcommand{\R}{\mathbb{R}}
\newcommand{\N}{\mathbb{N}}
\newcommand{\Pcal}{\mathcal{P}}
\newcommand{\Mcal}{\mathcal{M}}
\newcommand{\Qcal}{\mathcal{Q}}
\newcommand{\Zcal}{\mathcal{Z}}
\newcommand{\Fcal}{\mathcal{F}}
\newcommand{\Gcal}{\mathcal{G}}
\newcommand{\W}{\mathcal{W}}
\newcommand{\one}{\mathbf{1}}
\newcommand{\dd}{\,d}
\newcommand{\supp}{\operatorname{supp}}
\newcommand{\argmin}{\operatorname*{arg\,min}}
\newcommand{\argmax}{\operatorname*{arg\,max}}
\newcommand{\esssup}{\operatorname*{ess\,sup}}
\newcommand{\ri}{\operatorname{ri}}
\newcommand{\dom}{\operatorname{dom}}
\newcommand{\Law}{\operatorname{Law}}
\newcommand{\Var}{\operatorname{Var}}
\newcommand{\Cov}{\operatorname{Cov}}
\newcommand{\diag}{\operatorname{diag}}
\newcommand{\Id}{\mathrm{Id}}
\newcommand{\Lip}{\operatorname{Lip}}
\newcommand{\osc}{\operatorname{osc}}
\newcommand{\diam}{\operatorname{diam}}
\newcommand{\proj}{\operatorname{proj}}
\newcommand{\e}{\mathrm{e}}

\newenvironment{argument}{\paragraph{Argument.}\begin{enumerate}[leftmargin=2em]}{\end{enumerate}}
\newenvironment{proofoutline}{\paragraph{Proof architecture.}\begin{enumerate}[leftmargin=2em]}{\end{enumerate}}


% Source-faithfulness apparatus used by the paper reconstructions.
\newcommand{\sourceanchor}[1]{\par\smallskip\noindent\textbf{Source anchor.} #1\par\smallskip}
\newcommand{\sourcecomment}[1]{\begin{remark}#1\end{remark}}
\newenvironment{sourceguide}
  {\begin{description}[style=nextline,leftmargin=3.2cm,labelwidth=3cm]}
  {\end{description}}

% Backward-compatible environments used by some of the older reconstructions.
\newenvironment{distilled}[1][]{\begin{theorem}[#1]}{\end{theorem}}
\newenvironment{distillation}{\begin{enumerate}[leftmargin=2em]}{\end{enumerate}}

\newenvironment{commentarynotes}{\begin{remark}\begin{enumerate}[leftmargin=2em]}{\end{enumerate}\end{remark}}

\title{Modernized Proof Reconstruction of Generalized Schr\"odinger Bridge Matching}
\author{}
\date{Liu, Lipman, Nickel, Karrer, Theodorou, and Chen (2024)}
\begin{document}
\maketitle

\section*{Source map and scope}
\begin{sourceguide}
\item[Primary source] Liu et al., arXiv:2310.02233v2 (ICLR 2024 source PDF).
\item[Coverage] The mathematical development behind the GSBM alternating scheme: the fixed-marginal matching problem, conditional stochastic control, the quadratic closed form, path-integral reweighting, monotonicity, the fixed-point theorem, and the Appendix C derivations.
\item[Source anchors]
\begin{itemize}[leftmargin=2em]
  \item equations (1)--(11) and Section 3;
  \item Propositions 1--2, Lemma 3, Proposition 4, Theorems 5--6;
  \item Appendix B, equations (17)--(28);
  \item Appendix C.1, Lemma 7 and equations (29)--(30);
  \item Appendix C.2, equations (31)--(33);
  \item Appendix C.3, equations (34)--(35).
\end{itemize}
\end{sourceguide}

The exposition follows the source order and keeps the paper's proof devices.  In
particular, the conditional-control decomposition is proved by disintegration and
Fubini, the quadratic bridge is obtained from a cited killed-diffusion formula and
linear SDE calculations, and the path-integral formula is obtained by changing the
reference path law with Girsanov.  Later mathematical sources used to make omitted
steps explicit are cited at the point of use.

\section{The generalized Schr\"odinger bridge problem}
Let $(X_t)_{0\le t\le1}$ solve
\begin{equation}
  dX_t=u_t(X_t)\,dt+\sigma\,dW_t,
  \qquad X_0\sim\mu,
  \qquad X_1\sim\nu,
  \tag{1}\label{eq:gsbm-sde}
\end{equation}
with $\sigma\ge0$.  Its time marginal density $p_t$ is required to satisfy
\begin{equation}
  \partial_t p_t(x)
  =-\nabla\!\cdot\bigl(u_t(x)p_t(x)\bigr)
   +\frac{\sigma^2}{2}\Delta p_t(x),
  \qquad p_0=\mu,
  \qquad p_1=\nu.
  \tag{2}\label{eq:gsbm-fpe}
\end{equation}
Given a state cost $V_t:\R^d\to\R$, the generalized Schr\"odinger bridge
problem is
\begin{equation}
  \inf_{(p,u)}
  \int_0^1\int_{\R^d}
  \left(\frac12\|u_t(x)\|^2+V_t(x)\right)p_t(x)\,dx\,dt
  \quad\text{subject to \eqref{eq:gsbm-fpe}.}
  \tag{3}\label{eq:gsbm-main}
\end{equation}
When $V_t\equiv0$, this is the usual kinetic-action formulation of the
Schr\"odinger bridge.  The new point is that the interior marginal path is no
longer prescribed: it must be optimized jointly with the drift.

\section{Matching a prescribed marginal path}
The paper first recalls two matching procedures for a fixed smooth positive
probability path $(p_t)_{0\le t\le1}$.

\subsection{Implicit matching}
For a scalar potential $s_t(x)$, define
\begin{equation}
\begin{split}
  \mathcal L_{\mathrm{implicit}}(s)
  :=&\;\E_{\mu}[s_0(X_0)]-\E_{\nu}[s_1(X_1)]\\
   &+\int_0^1\E_{p_t}\left[
      \partial_t s_t(X_t)
      +\frac12\|\nabla s_t(X_t)\|^2
      +\frac{\sigma^2}{2}\Delta s_t(X_t)
   \right]dt.
\end{split}
\tag{4}\label{eq:gsbm-implicit}
\end{equation}
Assume $u_t$ is any drift generating $p_t$.  Multiply the Fokker--Planck
equation by $s_t$, integrate in $x$ and $t$, and integrate by parts.  Boundary
terms in time give
\[
  \E_{\nu}[s_1]-\E_{\mu}[s_0]
  =\int_0^1\E_{p_t}\left[
    \partial_t s_t+u_t\cdot\nabla s_t
    +\frac{\sigma^2}{2}\Delta s_t
  \right]dt.
\]
Substituting into \eqref{eq:gsbm-implicit} yields
\[
  \mathcal L_{\mathrm{implicit}}(s)
  =\int_0^1\E_{p_t}\left[
    \frac12\|\nabla s_t\|^2-u_t\cdot\nabla s_t
  \right]dt.
\]
If $u_t^*=\nabla s_t^*$ is the gradient drift of least kinetic energy among all
drifts producing $p_t$, then the weighted Helmholtz decomposition gives
\[
  u_t=u_t^*+w_t,
  \qquad
  \nabla\!\cdot(p_t w_t)=0,
\]
and hence
\[
  \int p_t w_t\cdot\nabla s_t\,dx=0.
\]
Completing the square gives
\begin{equation}
  \mathcal L_{\mathrm{implicit}}(s)
  -\mathcal L_{\mathrm{implicit}}(s^*)
  =\frac12\int_0^1\E_{p_t}
    \|\nabla s_t-\nabla s_t^*\|^2dt.
  \tag{4a}\label{eq:gsbm-implicit-square}
\end{equation}
Thus the minimizer is the unique gradient representative.  This is the
variational fact used in Proposition 1.  The weighted decomposition is the
modern form of the argument; see Ambrosio--Gigli--Savar\'e
\cite{AmbrosioGigliSavare2008} and the action-matching discussion cited by the
paper \cite{NeklyudovEtAl2023}.

\subsection{Explicit matching}
Suppose the prescribed path is represented as a mixture of endpoint-conditioned
paths:
\[
  p_t(x)=\int p_{t\mid0,1}(x\mid x_0,x_1)
  \,p_{0,1}(dx_0,dx_1).
\]
Assume each conditional path is generated by
\[
  dX_t=u_t(X_t\mid x_0,x_1)dt+\sigma dW_t.
\]
The paper's explicit regression loss is
\begin{equation}
  \mathcal L_{\mathrm{explicit}}(\theta)
  :=\int_0^1\E_{p_{0,1}}
    \E_{p_{t\mid0,1}}
    \left[\frac12
      \|u_t^\theta(X_t)-u_t(X_t\mid x_0,x_1)\|^2
    \right]dt.
  \tag{5}\label{eq:gsbm-explicit}
\end{equation}
Its pointwise regression minimizer is the conditional expectation of the
conditional drift given $X_t=x$.  Appendix C.1 proves that this projected drift
has the same Fokker--Planck marginals; the full calculation is reconstructed in
Section \ref{sec:gsbm-lemma7} below.

\section{The alternating GSBM decomposition}
The algorithm alternates between two exact mathematical subproblems.

\begin{proposition}[Proposition 1: Stage 1]
Fix the marginal path $(p_t)$.  The unique Stage-1 minimizer of
\eqref{eq:gsbm-main} is the gradient drift $\nabla s_t^*$ returned by implicit
matching.
\end{proposition}

\begin{proof}
For fixed $p_t$, the state-cost term
\[
  \int_0^1\E_{p_t}[V_t(X_t)]dt
\]
is constant in $u$.  Stage 1 therefore minimizes only the kinetic term among
drifts satisfying the fixed Fokker--Planck equation.  Equation
\eqref{eq:gsbm-implicit-square} identifies its unique gradient minimizer.
This is exactly the one-line argument printed in Appendix B.
\end{proof}

\subsection{Stage 2 as conditional stochastic control}
Let $p_{0,1}^\theta$ be the endpoint coupling induced by the current drift.  Stage
2 fixes this coupling and varies the conditional interior laws.

\begin{proposition}[Proposition 2: conditional SOC]
Suppose
\[
  p_t(x)=\int p_{t\mid0,1}(x\mid x_0,x_1)
  \,p_{0,1}^\theta(dx_0,dx_1).
\]
With $p_{0,1}^\theta$ fixed, Stage 2 decomposes into independent endpoint-pinned
control problems
\begin{equation}
\begin{aligned}
  \inf_{u_{\cdot\mid0,1}}
  \quad&\int_0^1
  \E_{p_{t\mid0,1}}
  \left[\frac12\|u_t(X_t\mid x_0,x_1)\|^2+V_t(X_t)\right]dt,\\
  \text{subject to}\quad&
  dX_t=u_t(X_t\mid x_0,x_1)dt+\sigma dW_t,\\
  &X_0=x_0,\qquad X_1=x_1.
\end{aligned}
\tag{6}\label{eq:gsbm-cond-soc}
\end{equation}
\end{proposition}

\begin{proof}
Start from the objective in \eqref{eq:gsbm-main}.  By the law of total
expectation and Fubini,
\begin{align*}
 &\int_0^1\E_{p_t}
   \left[\frac12\|u_t\|^2+V_t\right]dt\\
 &=\int_0^1\E_{p_{0,1}^\theta}
   \E_{p_{t\mid0,1}}
   \left[\frac12\|u_t\|^2+V_t\right]dt\\
 &=\E_{p_{0,1}^\theta}
   \left[
     \int_0^1\E_{p_{t\mid0,1}}
     \left[\frac12\|u_t\|^2+V_t\right]dt
   \right].
\end{align*}
Because the mixing law is fixed, the outer expectation separates the endpoint
pairs.

The same disintegration is applied term by term to the Fokker--Planck equation.
Under the regularity assumptions invoked in the paper
\cite{Anderson1982,YongZhou1999}, differentiation and integration may be
interchanged:
\[
  \partial_t p_t
  =\E_{p_{0,1}^\theta}[\partial_t p_{t\mid0,1}],
\]
\[
  \nabla\!\cdot(u_t p_t)
  =\E_{p_{0,1}^\theta}
    [\nabla\!\cdot(u_t p_{t\mid0,1})],
\]
and
\[
  \Delta p_t
  =\E_{p_{0,1}^\theta}[\Delta p_{t\mid0,1}].
\]
The endpoint laws disintegrate as
\[
  p_0=\E_{p_{0,1}^\theta}[\delta_{x_0}],
  \qquad
  p_1=\E_{p_{0,1}^\theta}[\delta_{x_1}].
\]
Therefore each conditional component is required to satisfy
\[
  \partial_t p_{t\mid0,1}
  =-\nabla\!\cdot(u_t p_{t\mid0,1})
   +\frac{\sigma^2}{2}\Delta p_{t\mid0,1},
  \qquad
  p_{0\mid0,1}=\delta_{x_0},
  \quad
  p_{1\mid0,1}=\delta_{x_1},
\]
which is the Fokker--Planck form of \eqref{eq:gsbm-cond-soc}.
\end{proof}

\begin{remark}[Coefficient dropped in the printed Appendix B line]
After correctly decomposing equation (17b), the paper's final displayed
conditional Fokker--Planck equation prints $+\Delta p_{t\mid0,1}$ rather than
$+(\sigma^2/2)\Delta p_{t\mid0,1}$.  The SDE in (6b), the preceding lines, and
all later formulas retain $\sigma$, so this is a typographical omission rather
than a change of model.
\end{remark}

\section{HJB and Hopf--Cole representation}
The paper states that the optimal conditional control has gradient form
\[
  u_t^*(x\mid x_0,x_1)=\sigma^2\nabla\log\Psi_t(x\mid x_1).
\]
Its displayed equation (18) is
\begin{equation}
  \partial_t\Psi_t(x\mid x_1)
  =-\frac{\sigma^2}{2}\Delta\Psi_t(x\mid x_1)
   +V_t(x)\Psi_t(x\mid x_1),
  \qquad
  \Psi_1(x\mid x_1)=\delta_{x_1}(x).
  \tag{18}\label{eq:gsbm-source-hopfcole}
\end{equation}
This is obtained from the HJB equation by an exponential transform.  With the
unscaled running cost in \eqref{eq:gsbm-cond-soc}, the standard transform
$\Psi=e^{-J/\sigma^2}$ produces a potential term $V_t/\sigma^2$ rather than
$V_t$.  Equation \eqref{eq:gsbm-source-hopfcole} is therefore consistent if
$V_t$ has already been rescaled by $\sigma^{-2}$; Proposition 4 later writes
that scaling explicitly.  The source does not state this convention at equation
(18), so both forms are recorded rather than silently identifying them.

\section{Quadratic state cost: Lemma 3}
Let
\[
  V(x)=\alpha\|\sigma x\|^2,
  \qquad
  \alpha>0,
  \quad \sigma>0,
  \qquad
  \eta:=\sigma\sqrt{2\alpha}.
\]

\begin{lemma}[Lemma 3]
The optimizer of \eqref{eq:gsbm-cond-soc} is Gaussian:
\[
  X_t^*\sim\mathcal N(c_t x_0+e_t x_1,\gamma_t^2 I_d),
\]
where
\begin{align}
  c_t&=\frac{\sinh(\eta(1-t))}{\sinh\eta},\label{eq:gsbm-ct}\\
  e_t&=\cosh(\eta(1-t))-c_t\cosh\eta,\label{eq:gsbm-et}\\
  \gamma_t&=\sigma
  \sqrt{\frac{\sinh(\eta(1-t))}{\eta}e_t}.
  \label{eq:gsbm-gammat}
\end{align}
As $\alpha\downarrow0$, these coefficients converge to those of the Brownian
bridge; sending $\sigma\downarrow0$ afterward gives the straight line.
\end{lemma}

\begin{proof}
The source imports the conditioned killed-diffusion formula of
Mazzolo--Monthus \cite[Eq.~(84)]{MazzoloMonthus2022}.  In the present notation,
the optimal drift is
\begin{equation}
  u_t(x\mid x_0,x_1)
  =\frac{\eta}{\sinh(\eta(1-t))}x_1
   -\frac{\eta}{\tanh(\eta(1-t))}x.
  \tag{19}\label{eq:gsbm-linear-drift}
\end{equation}
This is a linear SDE.  Its mean $\mu_t$ satisfies
\[
  \dot\mu_t
  =\frac{\eta}{\sinh(\eta(1-t))}x_1
   -\frac{\eta}{\tanh(\eta(1-t))}\mu_t.
\]
Let
\[
  g_t:=\exp\left(
    \int_0^t\frac{\eta}{\tanh(\eta(1-\tau))}\,d\tau
  \right).
\]
With $v_\tau=\sinh(\eta(1-\tau))$, the integral becomes a logarithmic
derivative, giving
\[
  g_t=\frac{\sinh\eta}{\sinh(\eta(1-t))}.
\]
Variation of constants yields the source's equation (20):
\[
  \mu_t
  =\frac{
      x_1\int_0^t
      \frac{\eta g_\tau}{\sinh(\eta(1-\tau))}\,d\tau+K
    }{g_t}.
\]
The initial condition $\mu_0=x_0$ gives $K=x_0$.  Evaluating the integral and
simplifying gives
\[
  \mu_t=c_t x_0+e_t x_1,
\]
with \eqref{eq:gsbm-ct}--\eqref{eq:gsbm-et}.

The covariance matrix $\Sigma_t$ obeys
\[
  \dot\Sigma_t
  =-\frac{2\eta}{\tanh(\eta(1-t))}\Sigma_t+\sigma^2 I_d,
  \qquad \Sigma_0=0.
\]
Applying the same integrating factor gives equation (22):
\begin{align*}
  \Sigma_t
  &=\frac{\sigma^2}{\eta}
  \sinh^2(\eta(1-t))
  \bigl(\coth(\eta(1-t))-\coth\eta\bigr)I_d\\
  &=\frac{\sigma^2}{\eta}
  \sinh(\eta(1-t))e_t I_d.
\end{align*}
Thus $\Sigma_t=\gamma_t^2I_d$ with \eqref{eq:gsbm-gammat}.

Finally, as $\eta\to0$,
\[
  c_t\to1-t,
  \qquad
  e_t\to t,
  \qquad
  \gamma_t\to\sigma\sqrt{t(1-t)},
\]
using the first-order expansions of $\sinh$ and $\cosh$.  These are the
Brownian bridge coefficients.  If $\sigma\to0$ as well, the variance vanishes
and the mean remains $(1-t)x_0+tx_1$.
\end{proof}

\section{Gaussian path approximation for general state costs}
For nonlinear $V_t$, the paper approximates each conditional optimizer by a
pinned isotropic Gaussian path
\begin{equation}
  p_t(\cdot\mid x_0,x_1)
  \approx\mathcal N(\mu_t,\gamma_t^2I_d),
  \quad
  \mu_0=x_0,\ \mu_1=x_1,
  \quad
  \gamma_0=\gamma_1=0.
  \tag{7}\label{eq:gsbm-gaussian-path}
\end{equation}
Writing
\[
  X_t=\mu_t+\gamma_t Z,
  \qquad Z\sim\mathcal N(0,I_d),
\]
gives the deterministic probability-flow velocity
\[
  \partial_t\mu_t+\frac{\dot\gamma_t}{\gamma_t}(X_t-\mu_t)
\]
and the score
\[
  \nabla\log p_t(X_t)=-\frac{1}{\gamma_t^2}(X_t-\mu_t).
\]
Adding $(\sigma^2/2)$ times the score produces the conditional SDE drift
\begin{equation}
  u_t(X_t\mid x_0,x_1)
  =\dot\mu_t+a_t(X_t-\mu_t),
  \qquad
  a_t:=\frac1{\gamma_t}
  \left(\dot\gamma_t-\frac{\sigma^2}{2\gamma_t}\right).
  \tag{8}\label{eq:gsbm-gaussian-drift}
\end{equation}
The paper parametrizes $\mu_t$ and $\gamma_t$ by splines and optimizes their
control points by Monte Carlo evaluation of the conditional objective.  This is
an approximation scheme, not an additional theorem about the exact optimizer.

\section{Path-integral reweighting: Proposition 4}
Let $q(d\bar X\mid x_0)$ denote Brownian path law with volatility $\sigma$ and
initial condition $x_0$.  Let $r(d\bar X\mid x_0,x_1)$ be a pinned proposal
associated with
\[
  dX_t=v_t(X_t)dt+\sigma dW_t,
  \qquad X_0=x_0,
  \quad X_1=x_1.
\]

\begin{proposition}[Proposition 4]
Assuming the absolute-continuity relation used in the source, the optimal path
law can be written
\begin{equation}
  p^*(d\bar X\mid x_0,x_1)
  =\frac1Z\omega(\bar X\mid x_0,x_1)
    r(d\bar X\mid x_0,x_1),
  \tag{10}\label{eq:gsbm-pi-law}
\end{equation}
where
\begin{equation}
\begin{split}
  \omega(\bar X\mid x_0,x_1)
  :=\exp\Bigg(&-\int_0^1
    \left[
      \frac{V_t(X_t)}{\sigma^2}
      +\frac{\|v_t(X_t)\|^2}{2\sigma^2}
    \right]dt\\
    &-\int_0^1\frac{v_t(X_t)^T}{\sigma}\,dW_t
  \Bigg).
\end{split}
\tag{11}\label{eq:gsbm-pi-weight}
\end{equation}
\end{proposition}

\begin{proof}
The source enforces the terminal pinning by adding the singular cost
$-\log\one_{\{x_1\}}(X_1)$ and dividing the whole objective by $\sigma^2$:
\begin{align}
  \inf_u\;&\int_0^1\E\left[
    \frac{\|u_t\|^2}{2\sigma^2}
    +\frac{V_t(X_t)}{\sigma^2}
  \right]dt
  -\E\log\one_{\{x_1\}}(X_1),
  \tag{23a}\label{eq:gsbm-singular-cost}\\
  &dX_t=u_t(X_t)dt+\sigma dW_t,
  \qquad X_0=x_0.
  \tag{23b}
\end{align}
The linearly solvable control identity of Kappen \cite{Kappen2005} gives
\begin{equation}
  p^*(d\bar X\mid x_0,x_1)
  =\frac1{Z'}
  \exp\left(-\int_0^1\frac{V_t(X_t)}{\sigma^2}dt\right)
  \one_{\{x_1\}}(X_1)
  q(d\bar X\mid x_0).
  \tag{24}\label{eq:gsbm-source-pi}
\end{equation}
To avoid sampling an endpoint event of negligible probability, the source
rebases \eqref{eq:gsbm-source-pi} to a pinned proposal $r$:
\begin{equation}
  p^*(d\bar X\mid x_0,x_1)
  =\frac1Z
  \exp\left(-\int_0^1\frac{V_t(X_t)}{\sigma^2}dt\right)
  \frac{dq}{dr}(\bar X)
  r(d\bar X\mid x_0,x_1).
  \tag{25}\label{eq:gsbm-rebase}
\end{equation}
The information-theoretic control references cited by the paper are
\cite{TheodorouBuchliSchaal2010,Theodorou2015}.  Girsanov's formula is then
written as
\begin{equation}
  \frac{dq}{dr}(\bar X)
  =\exp\left(
    -\int_0^1\frac{\|v_t(X_t)\|^2}{2\sigma^2}dt
    -\int_0^1\frac{v_t(X_t)^T}{\sigma}\,dW_t
  \right).
  \tag{26}\label{eq:gsbm-girsanov}
\end{equation}
Substituting \eqref{eq:gsbm-girsanov} into \eqref{eq:gsbm-rebase} gives
\eqref{eq:gsbm-pi-law}--\eqref{eq:gsbm-pi-weight}.
\end{proof}

\begin{remark}[Pinned-path singularity in the source formulation]
For a nondegenerate diffusion on $\R^d$, the event $X_1=x_1$ has probability
zero under unconditioned Brownian path law.  Consequently
$\one_{\{x_1\}}(X_1)$ and the assertion that an exactly pinned proposal is
absolutely continuous with respect to unconditioned $q$ are formal shorthand.
A rigorous version uses a Brownian bridge reference, a regular conditional law,
or an approximating terminal density.  The reconstruction retains the source's
indicator and Radon--Nikodym steps because those are the stated proof route,
while making the singularity explicit.
\end{remark}

\section{Monotonicity of exact alternating optimization}
Let $(p_t^{\theta_n},u_t^{\theta_n})$ be the result after $n$ iterations.  Write
$p_{0,1}^{\theta_n}$ for its endpoint coupling and
$p_{t\mid0,1}^{\theta_n}$ for its conditional law.  Let
$p_{t\mid0,1}^*$ be the exact minimizer of the conditional problem.  Stage 2
sets
\begin{equation}
  p_t^{\theta_{n+1}}
  :=\int p_{t\mid0,1}^*(\cdot\mid x_0,x_1)
  \,p_{0,1}^{\theta_n}(dx_0,dx_1).
  \tag{27}\label{eq:gsbm-next-marginal}
\end{equation}

\begin{theorem}[Theorem 5: local descent]
For exact Stage-1 and Stage-2 minimization,
\[
  \mathcal L(\theta_n)\ge\mathcal L(\theta_{n+1}).
\]
\end{theorem}

\begin{proof}
Disintegrate the $n$th objective:
\begin{align*}
  \mathcal L(\theta_n)
  &=\int p_{0,1}^{\theta_n}(dx_0,dx_1)
    \int_0^1\int
    \left[
      \frac12\|u_t^{\theta_n}\|^2+V_t
    \right]
    p_{t\mid0,1}^{\theta_n}(dx)dt.
\end{align*}
Conditional optimality in Stage 2 gives
\begin{align*}
  \mathcal L(\theta_n)
  &\ge\int p_{0,1}^{\theta_n}(dx_0,dx_1)
    \int_0^1\int
    \left[
      \frac12\|u_t^{\theta_n}\|^2+V_t
    \right]
    p_{t\mid0,1}^{*}(dx)dt\\
  &=\int_0^1\int
    \left[
      \frac12\|u_t^{\theta_n}(x)\|^2+V_t(x)
    \right]
    p_t^{\theta_{n+1}}(x)\,dxdt,
\end{align*}
where the last line uses \eqref{eq:gsbm-next-marginal}.  Stage 1 then replaces
$u^{\theta_n}$ by the minimum-kinetic-energy drift producing the new marginal
path.  Since $V_t$ depends only on that path,
\[
  \int\frac12\|u_t^{\theta_n}\|^2p_t^{\theta_{n+1}}
  \ge
  \int\frac12\|u_t^{\theta_{n+1}}\|^2p_t^{\theta_{n+1}}.
\]
Adding the unchanged state cost and integrating gives
\[
  \mathcal L(\theta_n)\ge\mathcal L(\theta_{n+1}).
\]
This is the chain printed as equation (28).
\end{proof}

\begin{remark}[Meaning of ``local convergence'']
The theorem proves monotonic nonincrease of objective values.  It does not by
itself prove convergence of the iterates, convergence to a stationary point,
or global optimality.  The source title ``Local convergence'' should therefore
be read as a descent property for exact alternating minimization.
\end{remark}

\section{The optimum is a fixed point}
\begin{theorem}[Theorem 6]
Every optimal solution $(p_t^*,u_t^*)$ of the generalized bridge problem is a
fixed point of the exact GSBM alternating scheme.
\end{theorem}

\begin{proof}
Disintegrate the optimal path law over its endpoint coupling:
\[
  p_t^*(x)=\int p_{t\mid0,1}^*(x\mid x_0,x_1)
  \,p_{0,1}^*(dx_0,dx_1).
\]
The proof has two source-stated steps.

\paragraph{Step 1: conditional optimality.}
The forward--backward SDE representation cited from Deep GSB
\cite{LiuEtAl2022DeepGSB} says that, for $p_{0,1}^*$-almost every endpoint
pair, the conditional optimal control satisfies the backward equation
associated with the HJB problem \eqref{eq:gsbm-source-hopfcole}.  The source
invokes uniqueness of that HJB/BSDE solution to conclude that
$p_{t\mid0,1}^*$ solves the conditional problem \eqref{eq:gsbm-cond-soc}.
Therefore Stage 2 leaves the optimal conditional mixture unchanged.

\paragraph{Step 2: matching returns the optimal Markov drift.}
The optimal drift $u_t^*$ is a gradient field in the cited GSB optimality
system.  Implicit matching returns the unique gradient field compatible with
$p_t^*$, so it returns $u_t^*$.  For explicit matching, the paper interprets
\eqref{eq:gsbm-explicit} as the Markovian projection of the reciprocal mixture
of endpoint bridges.  Since the mixture arising from the global optimum is
already the Markov law generated by $u_t^*$, its closest Markovian projection
is itself.  Hence explicit matching also returns $u_t^*$ at equilibrium.
\end{proof}

\begin{remark}[Dependence on external uniqueness results]
Theorem 6 is not proved solely from the algebra in the paper.  Its first step
uses the cited FBSDE representation and uniqueness of the conditional HJB/BSDE;
its second uses the Markovian-projection interpretation developed in the
Schr\"odinger bridge matching literature.  The reconstruction preserves this
citation-dependent proof rather than replacing it by a new variational
argument.
\end{remark}

\section{Appendix C.1: explicit matching preserves the marginals}
\label{sec:gsbm-lemma7}
\begin{lemma}[Lemma 7]
Let
\[
  p_t(x)=\E_{p_{0,1}}
  [p_{t\mid0,1}(x\mid X_0,X_1)]
\]
be a mixture of conditional diffusion paths.  Then the drift
\begin{equation}
  u_t^*(x)=
  \frac{
    \E_{p_{0,1}}
    [u_t(x\mid X_0,X_1)p_{t\mid0,1}(x\mid X_0,X_1)]
  }{p_t(x)}
  \tag{29}\label{eq:gsbm-markov-projection}
\end{equation}
produces the marginal path $p_t$.
\end{lemma}

\begin{proof}
Each conditional density satisfies
\[
  \partial_t p_{t\mid0,1}
  =-\nabla\!\cdot(u_{t\mid0,1}p_{t\mid0,1})
   +\frac{\sigma^2}{2}\Delta p_{t\mid0,1}.
\]
Average over $p_{0,1}$ and interchange expectation with differentiation:
\begin{align*}
  \partial_t p_t
  &=-\nabla\!\cdot
    \E_{p_{0,1}}[u_{t\mid0,1}p_{t\mid0,1}]
    +\frac{\sigma^2}{2}\Delta
    \E_{p_{0,1}}[p_{t\mid0,1}]\\
  &=-\nabla\!\cdot(u_t^*p_t)
    +\frac{\sigma^2}{2}\Delta p_t,
\end{align*}
where \eqref{eq:gsbm-markov-projection} is used in the last step.
\end{proof}

For any candidate $u_t^\theta(X_t)$, conditional least squares gives
\begin{align*}
 &\E\|u_t^\theta(X_t)-u_t(X_t\mid X_0,X_1)\|^2\\
 &=\E\|u_t^\theta(X_t)-u_t^*(X_t)\|^2
   +\E\|u_t^*(X_t)-u_t(X_t\mid X_0,X_1)\|^2.
\end{align*}
The second term is independent of $\theta$, so the paper writes
\begin{equation}
  \mathcal L_{\mathrm{explicit}}(\theta)
  =\E_{p_t}\left[
    \frac12\|u_t^\theta(X_t)-u_t^*(X_t)\|^2
  \right]+O(1).
  \tag{30}\label{eq:gsbm-explicit-square}
\end{equation}
Thus the explicit loss returns the Markovian projection
\eqref{eq:gsbm-markov-projection}.  The paper notes that this need not equal the
gradient minimizer of implicit matching away from equilibrium: the two differ by
a weighted divergence-free field.

\section{Appendix C.2: analytic drift and joint Gaussian sampling}
Starting from \eqref{eq:gsbm-gaussian-path}, the paper combines the
probability-flow velocity with half the diffusion score to obtain
\begin{equation}
  u_t(X_t\mid x_0,x_1)
  =\dot\mu_t
   +\frac1{\gamma_t}
    \left(\dot\gamma_t-\frac{\sigma^2}{2\gamma_t}\right)
    (X_t-\mu_t).
  \tag{31}\label{eq:gsbm-app-c-drift}
\end{equation}
This is \eqref{eq:gsbm-gaussian-drift}.

Write
\[
  a_t:=\frac1{\gamma_t}
  \left(\dot\gamma_t-\frac{\sigma^2}{2\gamma_t}\right),
  \qquad
  b_t:=\dot\mu_t-a_t\mu_t,
  \qquad
  g_t:=\int_0^t a_\tau d\tau.
\]
The linear SDE has the integrating-factor solution
\begin{equation}
  X_t=e^{g_t}
  \left(
    x_0+\int_0^t e^{-g_\tau}b_\tau d\tau
    +\int_0^t e^{-g_\tau}\sigma dW_\tau
  \right).
  \tag{32}\label{eq:gsbm-linear-solution}
\end{equation}
For $0\le s<t\le1$, independence of the It\^o increments on $[0,s]$ and
$[s,t]$, followed by It\^o isometry, gives
\begin{align*}
  \Cov(X_s,X_t)
  &=e^{g_s+g_t}
    \E\left[
      \left(\int_0^s e^{-g_\tau}\sigma dW_\tau\right)
      \left(\int_0^t e^{-g_\tau}\sigma dW_\tau\right)^T
    \right]\\
  &=e^{g_s+g_t}\int_0^s e^{-2g_\tau}\sigma^2d\tau\,I_d\\
  &=e^{g_t-g_s}\gamma_s^2I_d.
\end{align*}
By symmetry,
\begin{equation}
  \Cov(X_s,X_t)
  =\gamma_{\min(s,t)}^2
   e^{g_{\max(s,t)}-g_{\min(s,t)}}I_d.
  \tag{33}\label{eq:gsbm-covariance}
\end{equation}
This formula permits joint path samples at a time grid by one covariance matrix
and a Cholesky factorization rather than sequential Euler simulation.

\section{Appendix C.3: the deterministic quadratic case}
When $\sigma=0$, the conditional problem with quadratic state cost reduces to a
finite-horizon linear quadratic regulator:
\begin{equation}
  \inf_u\left
  \{x_1^TFx_1+
  \int_0^1(x_t^TQx_t+u_t^TRu_t)dt\right\},
  \qquad
  \dot x_t=Ax_t+Bu_t.
  \tag{34}\label{eq:gsbm-lqr}
\end{equation}
The optimal feedback is $u_t^*=-R^{-1}B^TP_tx_t$, where
\[
  -\dot P_t=A^TP_t+P_tA-P_tBR^{-1}B^TP_t+Q,
  \qquad P_1=F.
\]
Writing $H_t=P_t^{-1}$ gives the equivalent Lyapunov/Riccati equation
\begin{equation}
  \dot H_t=AH_t+H_tA^T-BR^{-1}B^T+H_tQH_t,
  \qquad H_1=F^{-1}.
  \tag{35}\label{eq:gsbm-inverse-riccati}
\end{equation}
The endpoint constraint is represented by $F\to\infty$, hence $H_1=0$.  For
$A=0$, $B=I_d$, $R=\frac12I_d$, and $Q=\alpha I_d$, the scalarized equation
has a hyperbolic-tangent solution, matching the zero-noise limit of the
stochastic formula.  This terminal-penalty argument is the one cited to the
linear-bridge literature \cite{ChenGeorgiou2015LinearBridges}.

\section{Source-fidelity notes}
\begin{commentarynotes}
\item Proposition 1 is genuinely a one-line reduction once the fixed-marginal
matching theorem is accepted; the longer derivation above makes that imported
fact explicit.
\item Proposition 2's printed final conditional Fokker--Planck line drops the
factor $\sigma^2/2$, although the source derivation and SDE retain it.
\item Equation (18) is ambiguous by a factor $\sigma^{-2}$ in the state-cost
term.  Proposition 4 uses the scaled potential explicitly.
\item Lemma 3 is not derived from first principles in GSBM: equation (19) is
imported from Mazzolo--Monthus, after which the mean and covariance ODEs are
solved in the paper.
\item Proposition 4 treats exact endpoint pinning through a singleton indicator
and an absolute-continuity change from unconditioned Brownian motion.  In
continuous state space this is formal shorthand for bridge conditioning or a
terminal-density limit.
\item Theorem 5 proves descent of objective values, not convergence of the
iterates.
\item Theorem 6 depends on cited FBSDE/HJB uniqueness and Markovian-projection
results; those citations are part of the source proof.
\end{commentarynotes}

\section*{Checklist of reconstructed arguments}
\begin{itemize}[leftmargin=2em]
\item Equations (1)--(5) and both matching objectives.
\item Proposition 1 and the kinetic-energy reduction.
\item Proposition 2 with objective and Fokker--Planck disintegration.
\item Equation (18) and its scaling ambiguity.
\item Lemma 3, including equations (19)--(22) and all limiting calculations.
\item Proposition 4, including equations (23)--(26) and Girsanov reweighting.
\item Theorem 5 and the full inequality chain (27)--(28).
\item Theorem 6 with both citation-dependent fixed-point steps.
\item Lemma 7 and equations (29)--(30).
\item Gaussian drift and covariance equations (31)--(33).
\item The deterministic LQR equations (34)--(35).
\end{itemize}

\begin{thebibliography}{99}
\bibitem{ChenGeorgiouPavon2021} Yongxin Chen, Tryphon T. Georgiou, and Michele Pavon. \emph{Stochastic Control Liaisons: Richard Sinkhorn Meets Gaspard Monge on a Schr\"odinger Bridge}. SIAM Review, 63(2):249--313, 2021. doi:10.1137/20M1339982. arXiv:2005.10963.
\bibitem{Leonard2014Survey} Christian L\'eonard. \emph{A Survey of the Schr\"odinger Problem and Some of Its Connections with Optimal Transport}. Discrete and Continuous Dynamical Systems--A, 34(4):1533--1574, 2014. doi:10.3934/dcds.2014.34.1533. arXiv:1308.0215.
\bibitem{Fortet1940} Robert Fortet. \emph{Resolution d'un systeme d'equations de M. Schrodinger}. Journal de mathematiques pures et appliquees, 9e serie, 19(1--4):83--105, 1940.
\bibitem{SinkhornKnopp1967} Richard Sinkhorn and Paul Knopp. \emph{Concerning Nonnegative Matrices and Doubly Stochastic Matrices}. Pacific Journal of Mathematics, 21(2):343--348, 1967.
\bibitem{FranklinLorenz1989} Joel Franklin and Jens Lorenz. \emph{On the Scaling of Multidimensional Matrices}. Linear Algebra and its Applications, 114/115:717--735, 1989. doi:10.1016/0024-3795(89)90490-4.
\bibitem{Carroll2004} Joseph E. Carroll. \emph{Birkhoff's Contraction Coefficient}. Linear Algebra and its Applications, 389:227--234, 2004. doi:10.1016/j.laa.2004.02.039.
\bibitem{EvesonNussbaum1995} Simon P. Eveson and Roger D. Nussbaum. \emph{An Elementary Proof of the Birkhoff--Hopf Theorem}. Mathematical Proceedings of the Cambridge Philosophical Society, 117(1):31--55, 1995.
\bibitem{BlanchetMurthy2019} Jose Blanchet and Karthyek R. A. Murthy. \emph{Quantifying Distributional Model Risk via Optimal Transport}. Mathematics of Operations Research, 44(2):565--600, 2019. doi:10.1287/moor.2018.0936. arXiv:1604.01446.
\bibitem{GaoKleywegt2023} Rui Gao and Anton J. Kleywegt. \emph{Distributionally Robust Stochastic Optimization with Wasserstein Distance}. Mathematics of Operations Research, 48(2):603--655, 2023. doi:10.1287/moor.2022.1275. arXiv:1604.02199.
\bibitem{MohajerinEsfahaniKuhn2018} Peyman Mohajerin Esfahani and Daniel Kuhn. \emph{Data-Driven Distributionally Robust Optimization Using the Wasserstein Metric: Performance Guarantees and Tractable Reformulations}. Mathematical Programming, 171(1--2):115--166, 2018. doi:10.1007/s10107-017-1172-1. arXiv:1505.05116.
\bibitem{WangGaoXie2025} Jie Wang, Rui Gao, and Yao Xie. \emph{Sinkhorn Distributionally Robust Optimization}. Operations Research, 2025. doi:10.1287/opre.2023.0294. arXiv:2109.11926.
\bibitem{Girsanov1960} I. V. Girsanov. \emph{On Transforming a Certain Class of Stochastic Processes by Absolutely Continuous Substitution of Measures}. Theory of Probability and Its Applications, 5(3):285--301, 1960. doi:10.1137/1105027.
\bibitem{CameronMartin1945} R. H. Cameron and W. T. Martin. \emph{Transformations of Wiener Integrals under a General Class of Linear Transformations}. Transactions of the American Mathematical Society, 58(2):184--219, 1945.
\bibitem{Yeh1978} J. Yeh. \emph{Transformation of Conditional Wiener Integrals under Translation and the Cameron--Martin Translation Theorem}. Tohoku Mathematical Journal, 30(4):505--515, 1978.
\bibitem{Kakutani1948} Shizuo Kakutani. \emph{On Equivalence of Infinite Product Measures}. Annals of Mathematics, 49(1):214--224, 1948.
\bibitem{ShiDeBortoliCampbellDoucet2023} Yuyang Shi, Valentin De Bortoli, Andrew Campbell, and Arnaud Doucet. \emph{Diffusion Schr\"odinger Bridge Matching}. Advances in Neural Information Processing Systems, 2023. arXiv:2303.16852.
\bibitem{LiuLipmanNickelKarrerTheodorouChen2024} Guan-Horng Liu, Yaron Lipman, Maximilian Nickel, Brian Karrer, Evangelos A. Theodorou, and Ricky T. Q. Chen. \emph{Generalized Schr\"odinger Bridge Matching}. International Conference on Learning Representations, 2024. arXiv:2310.02233.
\bibitem{DomingoEnrichDrozdzalKarrerChen2025} Carles Domingo-Enrich, Michal Drozdzal, Brian Karrer, and Ricky T. Q. Chen. \emph{Adjoint Matching: Fine-Tuning Flow and Diffusion Generative Models with Memoryless Stochastic Optimal Control}. arXiv preprint arXiv:2409.08861, 2025.

\bibitem{Birkhoff1957} Garrett Birkhoff. \emph{Extensions of Jentzsch's Theorem}. Transactions of the American Mathematical Society, 85(1):219--227, 1957.
\bibitem{Sinkhorn1964} Richard Sinkhorn. \emph{A Relationship Between Arbitrary Positive Matrices and Doubly Stochastic Matrices}. Annals of Mathematical Statistics, 35(2):876--879, 1964.
\bibitem{Sinkhorn1967Prescribed} Richard Sinkhorn. \emph{Diagonal Equivalence to Matrices with Prescribed Row and Column Sums}. American Mathematical Monthly, 74(4):402--405, 1967.
\bibitem{Bauer1965} Friedrich L. Bauer. \emph{An Elementary Proof of the Hopf Inequality for Positive Operators}. Numerische Mathematik, 7:331--337, 1965.
\bibitem{Bushell1973} P. J. Bushell. \emph{Hilbert's Metric and Positive Contraction Mappings in a Banach Space}. Archive for Rational Mechanics and Analysis, 52:330--338, 1973.

\end{thebibliography}

\end{document}
